\begin{table}
\centering
\begin{threeparttable}[!h]
\centering
\caption{\label{tab:partial_identification_region_coverage_right_tail_coarse_mild} Coverage of the simultaneous confidence region under mildly right-tail-coarse grids}
\centering
\begin{tabular}[t]{cccccccccc}
\toprule
\multicolumn{1}{c}{ } & \multicolumn{3}{c}{Coverage} & \multicolumn{3}{c}{Length of $\beta_0$ interval} & \multicolumn{3}{c}{Length of $\beta_1$ interval} \\
\cmidrule(l{3pt}r{3pt}){2-4} \cmidrule(l{3pt}r{3pt}){5-7} \cmidrule(l{3pt}r{3pt}){8-10}
$N$ & $K = 3$ & $K = 9$ & $K = 19$ & $K = 3$ & $K = 9$ & $K = 19$ & $K = 3$ & $K = 9$ & $K = 19$\\
\midrule
1000 & 0.9459 & 0.9389 & 0.9361 & 0.7577 & 0.4188 & 0.3023 & 0.9326 & 0.5498 & 0.4136\\
5000 & 0.9552 & 0.9474 & 0.9465 & 0.6554 & 0.3167 & 0.2009 & 0.7728 & 0.3930 & 0.2611\\
10000 & 0.9566 & 0.9553 & 0.9551 & 0.6298 & 0.2921 & 0.1765 & 0.7331 & 0.3548 & 0.2241\\
\bottomrule
\end{tabular}
\begin{tablenotes}[para,flushleft]
\footnotesize
\item \parbox{\linewidth}{Notes: Coverage reports the empirical coverage probability of the simultaneous confidence region for the true identified set. The length columns report the average lengths of the projection intervals.}
\end{tablenotes}
\end{threeparttable}
\end{table}